\begin{table}[htbp]
\centering
\caption{Item Completion under Open Tenders (18-month window)}
\label{tab:completion}
\begin{adjustbox}{max width=\textwidth}
\begin{threeparttable}
\small
\begin{tabular}{lcc}
\toprule
 & (1) Baseline & (2) + PBU FE \\
\midrule
$g65 \times Pre$ & 0.0969*** & 0.1017*** \\
 & (0.0037) & (0.0037) \\
\midrule
Observations & 878,888 & 878,870 \\
Item + Month FE & YES & YES \\
PBU FE & NO & YES \\
\bottomrule
\end{tabular}
\begin{tablenotes}
\small
\item \textit{Notes:} Dependent variable is an indicator equal to 1 if the
item was transacted (oc\_item\_status = 1) and 0 if no winner was selected.
Full 18-month sample of all purchase-offer-items (completed and not),
with item, month, and optional PBU fixed effects, sealed-bid and
log-quantity controls. Standard errors clustered at the item level.
Under open tenders, item completion is approximately 10 percentage
points higher (relative to a pre-period average of 66--77\%). The
completion effect matters for the fiscal-cost interpretation: the
headline price coefficient is estimated on completed items only,
whereas the static extrapolation applies only to that completed
subset. A structural counterfactual would need to price the shift in
who reaches completion.
*** \textit{p}$<$0.01, ** \textit{p}$<$0.05, * \textit{p}$<$0.1.
\end{tablenotes}
\end{threeparttable}
\end{adjustbox}
\end{table}
